% Mathematical Constants
% appendix_constants.tex

\section{Fundamental Constants}

\subsection{The Number $\pi$}
\[
\pi = 3.14159265358979323846264338327950288419716939937510\dots
\]
Euler computed $\pi$ to 150 decimal places in 1748 using the Gregory--Leibniz series:
\[
\frac{\pi}{4} = 1 - \frac{1}{3} + \frac{1}{5} - \frac{1}{7} + \frac{1}{9} - \cdots
\]
More rapidly converging formulas include Machin's identity:
\[
\frac{\pi}{4} = 4\arctan\frac{1}{5} - \arctan\frac{1}{239}
\]
Euler also showed that $\zeta(2) = \pi^2/6$, connecting $\pi$ to the sum of reciprocal squares.

\begin{lstlisting}[language=Python, caption=Computing $\pi$ using the Leibniz series]
from euler import leibniz_pi

# Slow convergence: needs millions of terms for 4 digits
pi_approx = leibniz_pi(1000000)
print(f"pi ≈ {pi_approx:.6f}")

# Better: use Euler's identity
import math
pi_euler = math.pi
print(f"math.pi = {pi_euler:.15f}")
\end{lstlisting}

\subsection{The Number $e$}
\[
e = 2.71828182845904523536028747135266249775724709369995\dots
\]
Euler proved $e$ is irrational in 1737. It arises naturally in many contexts:
\begin{itemize}
    \item Compound interest: $\displaystyle \lim_{n\to\infty}\left(1+\frac{1}{n}\right)^n = e$
    \item Series: $\displaystyle e = \sum_{n=0}^{\infty}\frac{1}{n!} = 1 + 1 + \frac{1}{2} + \frac{1}{6} + \cdots$
    \item Integral: $\displaystyle \int_1^e \frac{1}{x}\,dx = 1$
    \item Euler's identity: $e^{i\pi} + 1 = 0$
\end{itemize}

\begin{lstlisting}[language=Python, caption=Computing $e$ from its series]
from euler.infinite_series import series_for_e

e_approx = series_for_e(20)
print(f"e ≈ {e_approx:.15f}")

# Verify with math module
import math
print(f"math.e  = {math.e:.15f}")
\end{lstlisting}

\subsection{Euler's Constant $\gamma$}
\[
\gamma = 0.57721566490153286060651209008240243104215933593992\dots
\]
Also called the Euler--Mascheroni constant, it is defined as:
\[
\gamma = \lim_{n\to\infty}\left(H_n - \ln n\right)
\]
where $H_n = 1 + \frac{1}{2} + \frac{1}{3} + \cdots + \frac{1}{n}$ is the $n$-th harmonic number. Despite extensive study, it remains unknown whether $\gamma$ is rational or irrational.

\begin{lstlisting}[language=Python, caption=Computing $\gamma$]
from euler import euler_constant, euler_mascheroni_gamma

gamma1 = euler_constant(n=1000000)
gamma2 = euler_mascheroni_gamma()

print(f"gamma ≈ {gamma1:.15f}")
print(f"gamma ≈ {gamma2:.15f}")
\end{lstlisting}

\section{Algebraic Constants}

\subsection{The Golden Ratio $\phi$}
\[
\phi = \frac{1 + \sqrt{5}}{2} = 1.6180339887498948482045868343656381177203\dots
\]
The golden ratio satisfies $\phi^2 = \phi + 1$. It appears in the limit of consecutive Fibonacci numbers:
\[
\lim_{n\to\infty}\frac{F_{n+1}}{F_n} = \phi
\]
Its conjugate is $\psi = \frac{1-\sqrt{5}}{2} \approx -0.618$.

\subsection{The Silver Ratio $\delta_S$}
\[
\delta_S = 1 + \sqrt{2} = 2.4142135623730950488016887242096980785696\dots
\]
The silver ratio satisfies $\delta_S^2 = 2\delta_S + 1$. It appears in the geometry of the octagon and in the metallic means sequence.

\subsection{Square Roots}
\[
\sqrt{2} = 1.4142135623730950488016887242096980785696\dots
\]
The first known irrational number, discovered by the Pythagoreans. Babylonian tablets show it was known to 5 significant figures around 1800 BCE.

\[
\sqrt{3} = 1.7320508075688772935274463415058723669428\dots
\]
Appears in the altitude of an equilateral triangle of side length 2.

\section{Transcendental Constants}

\subsection{Natural Logarithms}
\[
\ln 2 = 0.6931471805599453094172321214581765680755\dots
\]
Appears in the alternating harmonic series: $\ln 2 = 1 - \frac{1}{2} + \frac{1}{3} - \frac{1}{4} + \cdots$

\[
\ln 10 = 2.3025850929940456840179914546843642076011\dots
\]
Converts between natural and common logarithms: $\log_{10} x = \frac{\ln x}{\ln 10}$.

\subsection{Catalan's Constant $G$}
\[
G = 0.9159655941772190150546035149323841107741\dots
\]
Defined by the series:
\[
G = \sum_{n=0}^{\infty}\frac{(-1)^n}{(2n+1)^2} = 1 - \frac{1}{9} + \frac{1}{25} - \frac{1}{49} + \cdots
\]
Appears in combinatorics, number theory, and the evaluation of certain integrals.

\subsection{Apéry's Constant $\zeta(3)$}
\[
\zeta(3) = 1.2020569031595942853997381615114499907649\dots
\]
Riemann's zeta function at $s=3$. Roger Apéry proved its irrationality in 1978, a landmark result. No simple closed form is known.

\subsection{Gompertz Constant $\delta$}
\[
\delta = 0.5963473623231940743410784993692793760741\dots
\]
Defined as $\delta = e \cdot \text{Ei}(-1)$, where $\text{Ei}$ is the exponential integral. Also equals $\int_0^{\infty}\frac{e^{-t}}{1+t}\,dt$.

\section{Constants from Number Theory}

\subsection{Khinchin's Constant $K$}
\[
K = 2.6854520010653064453097148354817956938201\dots
\]
For almost all real numbers, the geometric mean of the partial quotients in their continued fraction expansion equals $K$. Khinchin proved this in 1934.

\subsection{Feigenbaum's Constant $\delta$}
\[
\delta = 4.66920160910299067185320382\dots
\]
The universal ratio of successive bifurcation intervals in period-doubling routes to chaos. Discovered by Mitchell Feigenbaum in 1975, it appears in logistic maps and many other dynamical systems.

\section{Relations Between Constants}

Many constants are connected by elegant identities:

\begin{itemize}
    \item \textbf{Euler's identity:} $e^{i\pi} + 1 = 0$ connects $e$, $\pi$, $i$, $0$, and $1$.
    \item \textbf{Basel problem:} $\zeta(2) = \frac{\pi^2}{6}$.
    \item \textbf{Gamma function:} $\Gamma\left(\frac{1}{2}\right) = \sqrt{\pi}$.
    \item \textbf{Stirling's approximation:} $n! \sim \sqrt{2\pi n}\left(\frac{n}{e}\right)^n$.
    \item \textbf{Reflection formula:} $\Gamma(z)\Gamma(1-z) = \frac{\pi}{\sin(\pi z)}$.
    \item \textbf{Golden ratio and $\sqrt{5}$:} $\phi = \frac{1+\sqrt{5}}{2}$.
    \item \textbf{Logarithm of $e$:} $\ln e = 1$.
    \item \textbf{Riemann zeta at even integers:} $\zeta(2n) = \frac{(-1)^{n+1}B_{2n}(2\pi)^{2n}}{2(2n)!}$.
\end{itemize}

\section{Computing Constants}

All constants above can be computed in pure Python:

\begin{lstlisting}[language=Python, caption=Computing mathematical constants]
from euler import (
    euler_constant,        # gamma
    leibniz_pi,            # pi (slow)
    series_for_e,          # e
    zeta_function,         # zeta(s)
    zeta_even,             # zeta(2n)
)
from euler.special_functions import gamma_function
from math import pi, e, sqrt, log

# Fundamental constants
print(f"pi  = {pi:.15f}")
print(f"e   = {e:.15f}")
print(f"sqrt(2) = {sqrt(2):.15f}")
print(f"sqrt(3) = {sqrt(3):.15f}")

# Euler's constant
gamma = euler_constant(n=1000000)
print(f"gamma = {gamma:.15f}")

# Catalan's constant (via Dirichlet beta)
def catalan(n_terms=100000):
    s = 0.0
    for n in range(n_terms):
        s += (-1)**n / (2*n + 1)**2
    return s

G = catalan()
print(f"Catalan G = {G:.15f}")

# Apéry's constant zeta(3)
zeta3 = zeta_function(s=3.0, n_terms=1000000)
print(f"zeta(3) = {zeta3:.15f}")

# Verify Euler's identity
import cmath
lhs = cmath.exp(1j * pi) + 1
print(f"e^(i*pi) + 1 = {lhs}")  # ≈ 0
\end{lstlisting}

\begin{lstlisting}[language=Python, caption=Golden ratio and metallic means]
import math

# Golden ratio
phi = (1 + math.sqrt(5)) / 2
print(f"phi = {phi:.15f}")

# Silver ratio
delta_s = 1 + math.sqrt(2)
print(f"silver ratio = {delta_s:.15f}")

# Verify phi^2 = phi + 1
print(f"phi^2 = {phi**2:.15f}")
print(f"phi + 1 = {phi + 1:.15f}")

# Fibonacci limit
a, b = 1, 1
for _ in range(30):
    a, b = b, a + b
fib_limit = b / a
print(f"Fibonacci limit = {fib_limit:.15f}")
print(f"phi = {phi:.15f}")
\end{lstlisting}

\begin{exercise}
Compute $\ln 2$ using the alternating harmonic series to 10 decimal places. How many terms are needed?
\end{exercise}

\begin{exercise}
Prove that $\zeta(2) = \pi^2/6$ using Euler's product formula for $\sin x$.
\end{exercise}

\begin{exercise}
Compute Catalan's constant $G$ using the series definition. Compare your result with the known value $0.915965594\dots$
\end{exercise}
